\documentclass[10pt]{article}

\usepackage[margin=1in]{geometry}
\usepackage{amsmath}
\usepackage{booktabs}
\usepackage{enumitem}
\usepackage{hyperref}

\title{SEW-Offload: Static Expert-Window Scheduling for Prefetch-Hidden MoE Offloading on Ascend NPUs}
\author{}
\date{}

\begin{document}

\maketitle

\begin{abstract}
Mixture-of-Experts (MoE) inference stresses device memory because only a small
subset of experts is activated per token, while the full expert parameter set
still has to be stored somewhere. Existing MoE offloading systems are largely
designed around GPU assumptions: expert caching, host-to-device prefetching,
and overlap through CUDA streams. These assumptions become incomplete on Ascend
NPUs. Ascend execution is more sensitive to stable tensor addresses, fixed
shapes, graph replay, and explicitly scheduled data movement. This document
describes the design of SEW-Offload, a static expert-window scheduling system
for HBM-limited MoE inference on Ascend NPUs. SEW-Offload does not retrain the
router, does not alter top-k expert activation, and does not drop tokens.
Instead, it reinterprets expert offloading as a scheduling problem: fixed HBM
expert slots provide stable execution entries, deadline-aware prefetching
decides which missing experts should be loaded before they are needed, and
hit-first phased execution hides the remaining miss latency behind already
available expert computation. The goal is to minimize exposed prefetch stall,
not merely to maximize cache hit rate.
\end{abstract}

\section{Motivation}

Large MoE models such as Qwen3-30B-A3B activate only a small number of experts
for each token, but their total expert weights can still exceed the HBM budget
available to a single NPU card, especially when KV cache, long-context serving,
and concurrent requests compete for memory. Expert offloading is therefore a
natural way to serve larger MoE models under limited HBM. The central question
is not simply whether expert weights can be moved between host memory and HBM,
but whether the movement cost can be hidden well enough to preserve serving
latency and throughput.

GPU-oriented MoE offloading typically relies on a dynamic expert cache: keep
some experts resident, predict future experts, prefetch missing experts, and
overlap host-to-device copies with GPU computation. This abstraction is useful
but under-specifies the control surface exposed by Ascend NPUs. Ascend runtime
and kernels benefit from stable tensor addresses, fixed execution shapes, graph
capture and replay, and explicit memory movement orchestration. If an offloaded
expert is loaded into arbitrary weight tensors at runtime, the device may lose
the very regularity that makes Ascend execution efficient.

SEW-Offload starts from the observation that vLLM Ascend already provides a
grouped MoE execution representation: token-level routing results are converted
into per-expert token counts or group lists before grouped matrix
multiplication. Therefore, SEW-Offload must not claim ``per-token to
per-expert count conversion'' as its contribution. The missing piece is an
offloading-oriented execution window for expert weights: only a fixed number of
expert slots are present in HBM, each slot has a stable tensor address, and the
runtime dynamically maps active expert ids to those slots while orchestrating
loads, prefetches, replacements, and phased grouped execution.

\section{Problem Statement}

Given an MoE layer with a set of experts $E$, a limited HBM budget that can hold
only $S \ll |E|$ expert slots, and a stream of inference steps that produce
dynamic active expert sets, SEW-Offload aims to execute the original MoE
semantics while minimizing exposed offloading stall.

The system must preserve three correctness constraints:

\begin{itemize}[leftmargin=*]
  \item The router logits, top-k expert ids, and gate weights are unchanged.
  \item Every token is dispatched to the same experts as in the non-offloaded
  baseline.
  \item Expert outputs are combined with the same combine semantics as the
  existing vLLM Ascend grouped MoE path.
\end{itemize}

The performance objective is not raw cache hit rate. A miss is harmful only
when its loading time is exposed on the critical path. For a layer or execution
phase, the exposed stall can be written as:

\begin{equation}
T_{\text{stall}} =
\max(0, T_{\text{load-miss}} - T_{\text{overlap}}),
\end{equation}

where $T_{\text{load-miss}}$ is the host-to-HBM time required by missing
experts and $T_{\text{overlap}}$ is the amount of useful routing, dispatch,
resident-expert computation, later-layer computation, or cross-step computation
that can run while those experts are being loaded. SEW-Offload is designed to
increase $T_{\text{overlap}}$ and reduce the miss load that remains on the
critical path.

\section{Old Assumptions and New Control Surface}

Table~\ref{tab:assumptions} summarizes the design shift.

\begin{table}[h]
\centering
\small
\begin{tabular}{p{0.30\linewidth}p{0.30\linewidth}p{0.30\linewidth}}
\toprule
Aspect & GPU-style assumption & Ascend-aware redesign \\
\midrule
Expert residency &
Dynamic expert cache with flexible weight objects &
Fixed HBM expert slots with stable tensor addresses \\
Execution regularity &
Dynamic kernels and stream overlap are often acceptable &
Prefer fixed shapes, stable entries, graph replay, and few graph variants \\
Offloading objective &
Reduce miss count or copy time &
Reduce exposed stall after overlap \\
Miss handling &
Wait for missing experts, or prefetch predicted experts &
Execute hit experts first while miss experts are loaded \\
Dispatch representation &
Often discussed from token-level routing &
Reuse existing per-expert count/grouped execution; redesign weight residency \\
Data movement &
Mostly host-to-device copy scheduling &
Coordinate host-to-HBM loading with NPU-side movement and grouped compute \\
\bottomrule
\end{tabular}
\caption{SEW-Offload changes the optimization target from dynamic expert cache
management to static expert-window scheduling with prefetch-hidden execution.}
\label{tab:assumptions}
\end{table}

\section{Design Overview}

SEW-Offload consists of three coupled mechanisms.

\subsection{Static Expert Slots}

The HBM budget is represented as a fixed slot pool:

\[
\mathcal{S} = \{s_0, s_1, \ldots, s_{S-1}\}.
\]

Each slot owns preallocated weight tensors for the expert MLP weights. The
address, layout, dtype, and capacity of each slot are stable across inference
steps. The runtime maintains two mappings:

\[
\texttt{expert\_to\_slot}: E \rightarrow \mathcal{S} \cup \{\bot\},
\quad
\texttt{slot\_to\_expert}: \mathcal{S} \rightarrow E \cup \{\bot\}.
\]

Stable slots separate the dynamic expert identity from the static device-side
execution entry. The grouped MoE backend sees slot-local weights; the runtime
ensures that each slot contains the correct expert before it is used. This
design supports Ascend graph/static-kernel friendliness because the device-side
weight entry points are fixed even when the logical experts change.

\subsection{Deadline-Aware Prefetch and Orchestration}

The prefetch planner treats every candidate expert load as a job with a
deadline. For an expert $e$, the planner estimates:

\begin{itemize}[leftmargin=*]
  \item the next expected use time or layer/step deadline $d(e)$;
  \item the loading cost $c(e)$ from host memory into an HBM slot;
  \item the execution value $v(e)$, derived from expected token count, recent
  locality, and miss penalty;
  \item the replacement cost of evicting the current slot owner.
\end{itemize}

The planner prioritizes experts that are likely to be used soon, have high
token counts, and would cause large exposed stall if missing. Unlike a pure LRU
cache, the planner is allowed to choose a lower hit-rate schedule if that
schedule hides more loading time behind useful computation.

The most important prefetch opportunities are:

\begin{itemize}[leftmargin=*]
  \item \textbf{Intra-layer overlap:} load miss experts while already resident
  experts execute.
  \item \textbf{Inter-layer overlap:} use active experts from nearby layers and
  previous observations to prefetch before the next MoE layer is reached.
  \item \textbf{Cross-step decode overlap:} use decode-step locality to
  prefetch experts for the next token step while later layers of the current
  step are still executing.
\end{itemize}

\subsection{Hit-First Phased Execution}

When the current MoE layer produces an active expert set $A$, SEW-Offload
partitions it into resident hits and missing experts:

\[
A_{\text{hit}} = \{e \in A \mid \texttt{expert\_to\_slot}(e) \neq \bot\},
\quad
A_{\text{miss}} = A \setminus A_{\text{hit}}.
\]

Instead of blocking until all misses are loaded, SEW-Offload executes available
experts first. The layer is decomposed into a small number of grouped execution
phases:

\begin{enumerate}[leftmargin=*]
  \item immediately execute a grouped MLP phase for $A_{\text{hit}}$;
  \item concurrently load selected experts from $A_{\text{miss}}$ into
  available or replaceable slots;
  \item execute one or more follow-up grouped MLP phases as miss experts become
  ready;
  \item combine all partial expert outputs according to the original token and
  gate mapping.
\end{enumerate}

This mechanism converts a full-layer blocking miss into a phased execution
problem. It also avoids per-expert tiny kernels: ready experts are still batched
into grouped MLP phases, preserving the execution style of the existing Ascend
MoE backend.

\section{System Components}

\begin{table}[h]
\centering
\small
\begin{tabular}{p{0.24\linewidth}p{0.68\linewidth}}
\toprule
Component & Responsibility \\
\midrule
\texttt{ExpertStore} &
Owns the complete expert weights in host memory and exposes expert-level or
tile-level load interfaces. \\
\texttt{SlotPool} &
Preallocates fixed HBM expert slots and provides stable weight tensor entries. \\
\texttt{SlotManager} &
Maintains expert-to-slot mappings, handles admission, eviction, and slot
validity checks. \\
\texttt{WorkingSetObserver} &
Consumes existing vLLM Ascend routing outputs, especially per-expert token
counts and active expert ids, without changing router behavior. \\
\texttt{PrefetchPlanner} &
Scores candidate experts using deadline, token count, locality, and load cost;
emits asynchronous prefetch jobs. \\
\texttt{PhaseScheduler} &
Builds hit-first grouped execution phases and decides when to wait for
prefetched miss experts. \\
\texttt{ExecutionAdapter} &
Connects slot-local weights to the existing grouped MoE execution boundary with
minimal intrusion. \\
\texttt{Telemetry} &
Records hit rate, miss load time, exposed stall, copy/compute overlap, slot
churn, and correctness checks. \\
\bottomrule
\end{tabular}
\caption{SEW-Offload runtime components.}
\label{tab:components}
\end{table}

\section{Scheduling Algorithm}

At each MoE layer, the runtime performs the following steps:

\begin{verbatim}
Input:
  active experts A from existing grouped MoE metadata
  per-expert token counts C
  fixed slot pool S
  pending prefetch jobs P

1. Refresh slot readiness for completed prefetch jobs.
2. Partition A into hit experts A_hit and miss experts A_miss.
3. Launch or continue prefetch jobs for high-value miss experts.
4. Build phase 0 from A_hit and execute it immediately.
5. While unresolved experts remain:
     a. admit newly ready experts into a follow-up phase;
     b. if no ready expert exists, wait for the earliest useful prefetch;
     c. execute the ready phase through grouped MLP;
6. Combine phase outputs according to original routing and gate weights.
7. Use current observations to plan prefetches for future layers or steps.
\end{verbatim}

The planner uses a score such as:

\begin{equation}
\text{score}(e) =
\alpha \cdot \widehat{C}(e)
+ \beta \cdot \text{locality}(e)
+ \gamma \cdot \text{deadline\_urgency}(e)
- \delta \cdot \text{load\_cost}(e)
- \eta \cdot \text{eviction\_cost}(e),
\end{equation}

where $\widehat{C}(e)$ is the predicted or observed token count for expert
$e$. The exact score is a policy choice. The key point is that admission is
not optimized for hit rate alone; it is optimized for reducing the stall that
survives overlap.

\section{Ascend-Specific Design Rationale}

SEW-Offload is designed around Ascend-specific execution properties rather than
treating the NPU as a GPU replacement.

\textbf{Stable weight addresses.}
Fixed slots allow graph capture and static-kernel paths to observe stable
device-side weight entries. Logical expert ids may change, but the physical
slot tensors do not.

\textbf{Fixed execution windows.}
The number of HBM slots is fixed, and the runtime can restrict execution to a
small set of slot-window variants. This reduces dynamic graph diversity
compared with arbitrary expert weight objects.

\textbf{Explicit movement orchestration.}
Ascend systems expose data movement as a first-class optimization surface.
SEW-Offload separates host-to-HBM expert loading from NPU-side grouped compute
and records whether the copy time is hidden or exposed.

\textbf{Avoiding small expert kernels.}
Hit-first execution does not mean one kernel per expert. It groups ready
experts into phases so that the backend can continue using grouped matrix
multiplication.

\section{Integration with vLLM Ascend}

The intended implementation is low-intrusion. SEW-Offload should live in a new
package, for example \texttt{vllm\_ascend/moe\_offload/}, and remain disabled by
default. The first integration point is the expert execution boundary in the
existing fused MoE path, after routing and grouped metadata are available and
before the grouped expert MLP consumes weights.

The initial MVPs should be staged as follows:

\begin{enumerate}[leftmargin=*]
  \item \textbf{Observer mode:} record active experts, per-expert token counts,
  layer locality, and estimated reuse without changing execution.
  \item \textbf{Synchronous slot mode:} load missing experts into fixed slots
  synchronously to validate correctness under artificial slot budgets.
  \item \textbf{Prefetch mode:} add asynchronous prefetch jobs and measure
  exposed stall reduction.
  \item \textbf{Hit-first phased mode:} split layer execution into resident and
  ready-miss phases, while preserving grouped MLP execution.
  \item \textbf{Graph/static-kernel mode:} evaluate whether fixed slot windows
  improve graph replay stability and reduce overhead under selected capacity
  tiers.
\end{enumerate}

\section{Correctness and Failure Handling}

Correctness is protected by treating offloading as a weight-residency
transformation rather than a routing transformation. The runtime must verify
that every slot used in a phase contains the expected expert. A practical
implementation should maintain expert-id checksums or version counters for
slots, assert that phase-local expert ids match the original routed ids, and
compare per-layer outputs against a non-offloaded baseline during testing.

Failure handling should be conservative. If prefetch fails or the planner
cannot admit an expert before its deadline, the runtime falls back to a
synchronous load for that expert. This may hurt latency but preserves semantics.
If slot pressure is too high, the system should expose the stall in telemetry
rather than silently dropping experts or changing router behavior.

\section{Evaluation Plan}

The evaluation should answer four questions.

\begin{enumerate}[leftmargin=*]
  \item \textbf{Does fixed slotting preserve correctness?} Compare layer outputs
  and end-to-end generated tokens against the non-offloaded baseline.
  \item \textbf{Does orchestration reduce exposed stall?} Measure total
  host-to-HBM load time, hidden load time, and exposed stall separately.
  \item \textbf{Does hit-first phased execution improve latency under misses?}
  Compare blocking miss execution with phased execution under the same slot
  budget and workload.
  \item \textbf{Does Ascend-specific slot stability help graph/static execution?}
  Compare dynamic expert cache, fixed slots without graph replay, and fixed
  slots with selected graph/static-kernel variants.
\end{enumerate}

Key metrics include time-to-first-token, time-per-output-token, throughput,
P50/P95/P99 latency, slot hit rate, exposed miss stall, copy/compute overlap,
MTE/Cube overlap where available, slot churn, and HBM savings.

\section{Expected Contributions}

SEW-Offload is expected to contribute:

\begin{itemize}[leftmargin=*]
  \item an Ascend-aware reformulation of MoE offloading as static expert-window
  scheduling rather than dynamic expert caching alone;
  \item a fixed expert-slot abstraction that preserves stable HBM weight
  entries under dynamic expert identities;
  \item a deadline-aware prefetch planner that optimizes exposed stall instead
  of cache hit rate alone;
  \item a hit-first phased execution strategy that hides miss loading behind
  resident expert computation while retaining grouped MLP execution;
  \item a low-intrusion integration plan for vLLM Ascend that reuses existing
  per-expert count/grouped MoE metadata and does not modify router semantics.
\end{itemize}

\section{Limitations}

The current design assumes that host-to-HBM expert loading can be initiated
asynchronously and that the runtime can safely bind loaded weights into fixed
HBM slots before grouped execution. The design also assumes useful locality
exists in decode or continuous batching workloads; prefill-heavy workloads with
weak expert reuse may expose more load time. Finally, graph/static-kernel
benefits must be validated empirically on Ascend hardware, because dynamic
token counts may still create graph variants even when weight slots are stable.

\end{document}
